% Part: history
% Chapter: biographies
% Section: georg-cantor

\documentclass[../../../include/open-logic-section]{subfiles}

\begin{document}

\olfileid[tr]{his}{bio}{can}

\olsection{Georg Cantor}

\olphoto{cantor-georg}{Georg Cantor}

Georg Cantor'un (\textsc{gey}-org \textsc{kan}-tor) erken bir yaşam
öyküsü, onun Rusya'nın Sankt-Peterburg kentine giden bir gemide doğup
bulunduğunu ve anne babasının bilinmediğini ileri sürmüştür. Bu doğru
değildir; ancak Cantor gerçekten 1845'te Sankt-Peterburg'da doğmuştur.

Cantor matematik doktorasını 1867'de Berlin Üniversitesi'nden aldı. Küme
teorisi çalışmalarıyla tanınır ve küme teorisini ayrı bir araştırma alanı
olarak kuran kişi sayılır. Farklı büyüklüklerde sonsuz kümeler bulunduğunu
ilk ispatlayan odur. Teorileri, özellikle sonsuzluklar teorisi, dönemin
matematikçileri arasında yoğun tartışmalara yol açmış ve çalışmaları
tartışmalı bulunmuştur.

Cantor'un dinsel inançlarıyla matematik çalışmaları ayrılmaz biçimde
bağlıydı; sonlu ötesi sayılar teorisinin kendisine doğrudan Tanrı tarafından
iletildiğini bile ileri sürdü. Cantor yaşamının ilerleyen döneminde ruhsal
hastalık yaşadı. 1894'ten başlayarak, özellikle son yıllarında giderek daha
sık hastaneye yatırıldı. Matematikçi Leopold Kronecker'le arasının açılması
da dâhil olmak üzere çalışmalarına yöneltilen ağır eleştiriler depresyona
ve matematiğe ilgisinin azalmasına yol açtı. Depresyon dönemlerinde
felsefeye ve edebiyata yönelen Cantor, Shakespeare'in oyunlarını Francis
Bacon'ın yazdığına ilişkin bir teori bile yayımladı.

Cantor, 6 Ocak 1918'de Halle'deki bir sanatoryumda öldü.

\begin{reading}
Cantor'un kapsamlı yaşam öyküleri için \citet{Dauben1990} ve
\citet{Grattan-Guinness1971} kaynaklarına bakın. Cantor'un radikal
görüşleri BBC Radio 4'ün \emph{A Brief History of Mathematics} programında
da anlatılır \citep{Sautoy2014}. Cantor'un teorilerini rap biçiminde
dinlemek isterseniz \citet{Rose2012} kaynağına bakın.
\end{reading}

\end{document}
